(1) 介质静止系顺向光速 \(u'=c/n\)，实验室系速度为 \(u_+=(u'+v)/(1+u'v/c^2)\)。展开 \(1/(1+v/(nc))\) 一阶得 \(u_+=c/n+v(1-1/n^2)+O(v^2/c)\)。逆向波在介质系速度 \(-c/n\)，实验室系相对于逆方向的正速率为 \(u_-=c/n-v(1-1/n^2)+O(v^2/c)\)。系数 \(1-1/n^2\) 即 Fresnel 拖曳因子。

(2) 经过实验室固定长度 \(L\) 的传播时间分别为 \(t_+=L/u_+\)、\(t_-=L/u_-\)。设 \(a=c/n,\ \delta=v(1-1/n^2)\)，则 \(t_--t_+=L[(a-\delta)^{-1}-(a+\delta)^{-1}]\simeq2L\delta/a^2\)。故 \(\Delta t\simeq2Lv(n^2-1)/c^2\)。这里 \(L\) 按每束光实际穿过的介质长度定义；若装置存在多段反向流动，须带符号逐段相加而非额外凭空乘2。

(3) 光源实验室角频率 \(\omega_0=2\pi c/\lambda\)，假设两路合束时频率相同。相位差 \(\Delta\phi=\omega_0\Delta t\simeq4\pi Lv(n^2-1)/(\lambda c)\)，对应条纹位移数 \(\Delta N=\Delta\phi/(2\pi)=2Lv(n^2-1)/(\lambda c)\)。运动方向反转使相位差变号，实际比较正、反流时的条纹差应对这两个相反相位再取差，不能把它与单次相位变化混淆。

(4) 若 \(n\to1\)，速度合成给 \(u_\pm\to c\)，所以相位差消失，符合真空光速不变量。若 \(v\to0\)，两路时间一致且 \(\Delta\phi\to0\)。对强色散介质，介质系的相位传播与群传播以及折射率的频率依赖必须区分，单一 \(n\) 的 Fresnel 系数不再给出任意光脉冲的群时延；本题明确限制于无色散单频波与一阶速度展开。
